% Chapter VI detector-to-record successor, 17 October 2026.
% LaTeX fragment; requires amsmath and amssymb, no preamble.
% Agent-originated receiving construction; no thermal rate or native run assigned.

\section{A finite pointer-to-memory instrument with causal local access}
\label{sec:chvi-memory-channel}

\paragraph{Source and scope.}
The retained SAM\_land donor readout construction supplies neutral record
bits with degenerate rest channels and positive bounded couplings proportional
to $I-X$; its finite mark/write/unmark law retains the complete complementary
cloud. Its final record is held, with reset treated separately
(\path{samland_donor_readout1/DERIVATION.md}, Section~4, and its
\path{ENERGY_ACCOUNT.md}). We reuse that kind of record recipient below,
with a new, explicitly specified smooth coordinate response. The physical
memory covariance of \path{volii_physical_memory1} and the colored
classical field bath of \path{voli_radiative_memory1} retain their own
source laws; neither supplies a spin-register flip rate here.

The new instrument, finite-write estimate and correlated-hold bounds are
agent-originated derivations by OpenAI ChatGPT and Codex. They join a
decoded bit to the actual detector pointer under the stated receiving law.
The coordinate-to-qubit coupling specified below is an explicit apparatus
assignment; it is not derived here from a microscopic material actuator.
They do not extract a new positive cap mean from the old field, assign a
microscopic material or temperature, or replace the later two-arm transport
calculation. Local controls in this section require the decoded bit, detector
pointer and memory to be physically co-located, or a previously specified
causal delivery channel bringing the bit to that location.

\subsection{A correlated loaded pointer and its complete binary conversion}

Let $B$ be the isolated decoder-output qubit from the source-only protocol,
with $P_b=|b\rangle\langle b|$, $b\in\{0,1\}$. For each operational input
label $j$, let $\rho_j$ be the finite-energy reference output of the actual
decoder pulses on the ideal prepared marker. Its error obeys
$\operatorname{Tr}(P_{1-j}\rho_j)\leq e_{{\rm dec,ref},j}$, excluding
the marker preparation distance $\delta_{\rm prep}$. The field can remain
correlated with this bit. Let the actual analog detector noise observable be
$N$, acting on the pointer/source/field factors outside $B$, so
$[N,P_b]=0$. Assume the retained full detector calculation supplies the
per-label moment bound
\begin{equation}\begin{gathered}
\operatorname{Tr}(\rho_jN^2)\leq\mathcal B_{\rm ret}
 \\
\text{uniformly over the declared}\\\text{analog receiving interval.}
\label{eq:chvi-memory-noise-input}
\end{gathered}\end{equation}
The joint state $\rho_j$ may correlate $B$, $N$ and all other degrees of
freedom. No conditioning on an exceptionally rare bit value is used to
claim the same moment bound.

After the detector field gate, the specified bit-controlled translations
load the same pointer with signs $+a$ for $b=0$ and $-a$ for $b=1$.
Its loaded coordinate is therefore
\begin{equation}
\begin{gathered}
Q=N+a(P_0-P_1),\\
 a>0.
\end{gathered}
\label{eq:chvi-memory-loaded-coordinate}
\end{equation}
Equivalently, in a Schr\"odinger description the loading unitary translates
the ordinary pointer coordinate and the displayed equation specifies its
input Heisenberg observable. The source-defined field distribution and its
correlations remain in $N$. The added sign is the decoded bit's paid
apparatus displacement, not an assumed intrinsic mean of $N$.

Choose $0<w<a$. With $x=(q+w)/(2w)$ define the $C^1$ threshold response
\begin{equation}
\begin{gathered}
A_w(q)=
 \begin{cases}
 1,&q\leq-w,\\
 1-3x^2+2x^3,&-w<q<w,\\
 0,&q\geq w.
 \end{cases}
 \\
 0\leq A_w\leq1,
 \\
 G:=\|A_w'\|_\infty=\frac{3}{4w}.
\end{gathered}
\label{eq:chvi-memory-threshold}
\end{equation}
This $A_w$ is an apparatus response, not the accumulation scalar $A(r)$
or the computational \texttt{A\_OPERATOR}.

For a finite registration aperture, take $L_q>w$ and a $C^1$ multiplier
$\chi_{L_q}$ equal to one on $|q|\leq L_q$ and zero on
$|q|\geq2L_q$, with cubic shoulders and
$\|\chi_{L_q}'\|_\infty\leq3/(2L_q)$.
Use $A_{w,L_q}=\chi_{L_q}A_w$ in the physical memory pulse.
Its response remains in $[0,1]$ and has compact coordinate support.
The two derivative supports are disjoint, so
$\|A_{w,L_q}'\|_\infty\leq\max\{G,3/(2L_q)\}$.
The finite-pulse derivative bounds below apply with that derivative bound.
The uncut threshold formulas acquire the following additional wrong-bit
probability, charged once through the final aperture/read allowance:
\begin{equation}
 \epsilon_{\rm ap}\leq
 \frac{(\sqrt{\mathcal B_{\rm ret}}+a)^2}{L_q^2}
 \leq\frac{2(\mathcal B_{\rm ret}+a^2)}{L_q^2}.
 \label{eq:chvi-memory-aperture}
\end{equation}
This follows from $Q=N+aZ_B$ and the second-moment triangle inequality.
The trusted-band bound is supplemented by this aperture term; its output
is still part of the complete binary experiment. For $L_q=2^{20}$ and
$w=2^9$, the original $G$ is unchanged. The final budget can absorb this
probability into a stated combined aperture/read allowance.

Supply one memory qubit $R$, initially $|0\rangle_R$, with degenerate
rest Hamiltonian, and set $X_-=(I-X_R)/2$ and
$D_w=A_w(Q)\otimes X_-$. The ideal positive write pulse has area $\pi$:
\begin{equation}
\begin{gathered}
V_\theta=e^{-i\theta D_w},\quad 0\leq\theta\leq\pi,
 \\
 U_w=V_\pi.
\end{gathered}
\label{eq:chvi-memory-ideal-write}
\end{equation}
It is exactly the identity on the memory for $q\geq w$, and exactly
$X_R$ for $q\leq-w$. At every coordinate it is unitary, including the
finite transition band. Reading the memory in its $Z_R$ basis gives the
two Kraus operators on the retained incoming system
\begin{equation}
\begin{gathered}
K_0(Q)=e^{-i\pi A_w(Q)/2}\cos[\pi A_w(Q)/2],\\
 K_1(Q)=i e^{-i\pi A_w(Q)/2}\sin[\pi A_w(Q)/2],
\end{gathered}
\label{eq:chvi-memory-kraus}
\end{equation}
with $K_0^*K_0+K_1^*K_1=I$. Both outcomes and their conditional residual
source/pointer/field states are retained. The transition band is not
postselected away; its possibly erroneous binary output is included.
The memory need not be read immediately: Eq.~\eqref{eq:chvi-memory-ideal-write}
first creates a coherent memory correlation, to which the hold channel
below can be applied.

\noindent\textbf{Proposition (threshold error without independence).}
For each label $j$, the ideal write obeys
\begin{equation}
 e_{{\rm thresh,ref},j}
 \leq\min\left\{1,e_{{\rm dec,ref},j}
                  +\frac{\mathcal B_{\rm ret}}{(a-w)^2}\right\}.
\label{eq:chvi-memory-threshold-error}
\end{equation}
\emph{Proof.} On the simultaneous spectral subspace where the decoded
bit equals $j$ and $|N|<a-w$, Eq.~\eqref{eq:chvi-memory-loaded-coordinate}
lies strictly beyond the correct side of the transition band. The wrong
memory effect vanishes there. Since it is a positive contraction, the
commuting spectral calculus gives
\begin{equation}
 K_{1-j}(Q)^*K_{1-j}(Q)
 \leq P_{1-j}+\mathbf1_{\{|N|\geq a-w\}}.
\end{equation}
The spectral second-moment inequality bounds the second term's expectation
by $\mathcal B_{\rm ret}/(a-w)^2$. Taking the expectation in the full joint
state establishes Eq.~\eqref{eq:chvi-memory-threshold-error}, with no
factorization of bit and detector noise. $\square$

\subsection{A finite positive pulse with the pointer kinetic energy retained}

Let the pointer's actual rest Hamiltonian be $H_P=kP^2$, $k>0$,
$P=-i\partial_Q$, and retain it throughout the write. A rectangular pulse
of duration $\tau_w>0$ is
\begin{equation}
\begin{gathered}
H_{\rm write}=kP^2+\frac{\pi}{\tau_w}D_w,
 \\
 U_{\rm write}=e^{-i\tau_wH_{\rm write}}.
\end{gathered}
\label{eq:chvi-memory-finite-write}
\end{equation}
Both terms are positive; there is no kinetic parking. The actual
source--field generator $H_{\rm sf}(t)$ is a spectator commuting with these pointer
and memory operators. It factors from the comparison even when its state
is entangled with the pointer. Its own evolution and energy are retained.
Equation~\eqref{eq:chvi-memory-finite-write} specifies the write-stage
memory law: any additional noncommuting bath coupling during that stage
needs a separate write-error bound. Activating a different hold coupling
afterward retains its switching work rather than parking the pointer.

Suppose the joint write-entry state has $\langle P^2\rangle\leq p_2$.
In the intended pipeline the detector gate and bit-controlled translations
depend on the analog pointer momentum and commute with $P^2$; its free
analog hold does also. Thus the declared Gaussian input moment
$p_2=1/(4\sigma^2)$ travels unchanged through those stages. A further
control that does not preserve this moment needs its own bound.
Purification includes arbitrary correlations with the remaining systems.
The bounded $C^1$ multiplier $A_w$ preserves $D(P)$, and for a normalized
purification $\psi$,
\begin{equation}
\begin{gathered}
\|PV_\theta\psi\|\leq\sqrt{p_2}+|\theta|G,
 \\
 \|PD_wV_\theta\psi\|\leq\sqrt{p_2}+(|\theta|+1)G.
\end{gathered}
\label{eq:chvi-memory-derivative-estimates}
\end{equation}
These follow from $[P,D_w]=-iA_w'(Q)X_-$ and differentiation of the
commuting multiplication exponential. In particular both ideal orbits
needed by the bounded-control comparison have free-energy form expectation
at most
\begin{equation}
\begin{gathered}
C_w=k[\sqrt{p_2}+(\pi+1)G]^2
 \leq C_w^{\rm up}:=2k[p_2+(\pi_*+1)^2G^2],\\
\pi_*\geq\pi.
\end{gathered}
\label{eq:chvi-memory-write-energy}
\end{equation}
No $D(P^2)$ assumption on the initial state is required for this estimate.
The self-adjoint Hamiltonian in Eq.~\eqref{eq:chvi-memory-finite-write}
is the nonnegative free Hamiltonian plus a bounded self-adjoint control.

\noindent\textbf{Proposition (finite-write discrepancy).}
For the specified reachable input family,
\begin{equation}
\begin{aligned}
 \|(U_{\rm write}-U_w)\psi\|
 &\leq r_w:=(1+2\pi)\sqrt{2\tau_w C_w}\\
 &\leq(1+2\pi_*)\sqrt{2\tau_w C_w^{\rm up}}.
\end{aligned}
\label{eq:chvi-memory-write-discrepancy}
\end{equation}
The corresponding complete joint-output trace distance and each
memory-outcome probability error are at most $\min\{1,r_w\}$.

\emph{Proof.} Set $F(t)=e^{-itH_P}$ and move the actual pulse into its
free interaction picture. Its bounded generator is
$(\pi/\tau_w)F(t)^*D_wF(t)$; the ideal one is
$(\pi/\tau_w)D_w$. For any form-domain vector $\xi$, the spectral theorem
gives
\begin{equation}
 \|(F(t)-I)\xi\|^2\leq2|t|\langle\xi,H_P\xi\rangle.
\end{equation}
Expanding
$F(t)^*D_wF(t)-D_w
 =F(t)^*D_w(F(t)-I)+(F(t)^*-I)D_w$ and using
Eq.~\eqref{eq:chvi-memory-write-energy} bounds its action on the ideal
orbit by $2\sqrt{2tC_w}$. Bounded-generator Duhamel therefore bounds
the interaction-picture difference by
$2\pi\sqrt{2\tau_wC_w}$, using $t\leq\tau_w$ under the integral.
Returning from that picture costs at most
$\sqrt{2\tau_wC_w}$ on the final ideal orbit. Add the two bounds.
Purifying a mixed input and then taking partial traces gives the same
output trace-distance conclusion. $\square$

\subsection{Passive hold and correlated transverse disturbances}

After writing, turn off the write pulse. The memory rest law is the same
degenerate channel used by the retained SAM\_land record construction.
A more general admitted passive hold is
\begin{equation}
\begin{gathered}
H_d(t)=|0\rangle\langle0|_R\otimes L_0(t)
       +|1\rangle\langle1|_R\otimes L_1(t),
 \\
 [H_d(t),Z_R]=0,
\end{gathered}
\label{eq:chvi-memory-diagonal-hold}
\end{equation}
with well-defined block propagators on the retained environment. This
environment includes the pointer, source, field and every noise reference
needed to describe correlations. The memory may dephase or acquire joint
environmental records, but each $Z_R$ population is exactly constant.
Thus its readable binary record is held for every declared finite interval
under this law. No refresh or measurement is needed during the hold.

For an actual transverse disturbance, split any admitted additional
interaction into its diagonal part, absorbed in $H_d$, and its bounded
off-diagonal part
\begin{equation}
\begin{gathered}
V_{\rm off}(t)=\Pi_0V(t)\Pi_1+\Pi_1V(t)\Pi_0,
 \quad \Pi_b=|b\rangle\langle b|_R\otimes I,
 \\
 \Lambda_T=\int_0^T\|V_{\rm off}(u)\|\,\mathrm du.
\end{gathered}
\label{eq:chvi-memory-hold-exposure}
\end{equation}
The norm is on the full retained joint system, not an average of independent
cell errors. Interaction-picture bounded-perturbation Duhamel gives a
joint-channel trace-distance error at most $\min\{1,\Lambda_T\}$ on every
input and reference. Consequently the wrong-label population can increase
by at most $\min\{1,\Lambda_T\}$ uniformly for $0\leq t\leq T$.

A sharper population envelope is also available. If its initial value is
$e_0$, then
\begin{equation}
 e(t)\leq
 \sin^2\!\left(\min\left\{\frac\pi2,
                  \arcsin\sqrt{e_0}+\Lambda_t\right\}\right).
\label{eq:chvi-memory-population-envelope}
\end{equation}
To see this, purify the complete state. In the diagonal interaction picture,
the wrong-sector probability obeys, almost everywhere,
$|\dot e|\leq2\|V_{\rm off}\|\sqrt{e(1-e)}$ by Cauchy--Schwarz between
the two orthogonal sectors. Integrating the bound on
$\arcsin\sqrt e$ gives Eq.~\eqref{eq:chvi-memory-population-envelope},
with endpoint values obtained by continuity. For $e_0=0$ this becomes
$e(t)\leq\sin^2(\min\{\pi/2,\Lambda_t\})$.
Use this population envelope or the additive channel error in a budget;
do not charge the same disturbance twice.

This is a passive finite-spin receiving model, not an assigned Markov
flip rate. For example a finite correlated thermal environment may be
declared by a positive Ising Hamiltonian
\begin{equation}
\begin{gathered}
H_E=\sum_\ell h_\ell\frac{I-Z_\ell}{2}
       +\sum_\ell J_\ell\frac{I-Z_\ell Z_{\ell+1}}{2},
 \\
 h_\ell,J_\ell\geq0,
 \quad \rho_E=\frac{e^{-\beta H_E}}{\operatorname{Tr}e^{-\beta H_E}},
\end{gathered}
\end{equation}
and a static coupling $V_{\rm off}=X_R\otimes B_x+Y_R\otimes B_y$.
Its actual matrices give
$\Lambda_T\leq T(\|B_x\|+\|B_y\|)$, regardless of thermal correlations.
Diagonal $Z_R$ couplings may be retained in Eq.~\eqref{eq:chvi-memory-diagonal-hold}.
If a nonnegative finite total bath Hamiltonian is desired, include the
common offset $\|V_{\rm off}\|I$ (and the norm of any additional bounded
diagonal coupling) in its energy account. This does not change the memory
probabilities.
The bath preparation, $\beta$, couplings and energy offsets would be new
receiving operands; no such numerical thermal model is inferred from an
isotope covariance or inserted into the example below. No thermal activation
advantage is claimed by this state-independent norm bound.

\subsection{Readout, complete outcomes and causal access}

For local readout supply a blank qubit $A$. A positive bounded pulse
\begin{equation}
\begin{gathered}
H_{\rm copy}(t)=u(t)\Pi_1\otimes X_{-,A},\\
 u(t)\geq0,\quad\int u(t)\,\mathrm dt=\pi
\end{gathered}
\label{eq:chvi-memory-read-copy}
\end{equation}
copies the memory's orthogonal $Z_R$ label to $A$. It commutes with the
diagonal hold law, even with the full correlated environment retained.
The final $Z_A$ measurement uses the retained Born/instrument premise.
This specifies its ideal back-action: it preserves memory label populations,
while permitting correlation and dephasing. An actual copy/read channel
may instead carry a declared outcome error $\delta_{\rm read}$ uniform over
the receiving interval. A transverse exposure during copying is included
either in the extended hold exposure or in that read error, once.

More explicitly, first append the blank memory, apply the finite write,
then the full hold channel, append the read qubit and apply its copy unitary.
Projection onto $|b\rangle_A$ followed by the read-register partial trace
defines a completely positive instrument map $\mathcal J_b^t$ on the
retained systems. The two maps sum to a trace-preserving channel.
No residual source, pointer or environmental cloud is replaced by a pure
state merely because a classical outcome has been registered.
For an input already carrying the chapter's ready/failure flags, define
\begin{equation}
\begin{gathered}
\mathcal I_b^t(\widehat\rho)
     =\mathcal J_b^t(P_{\rm ready}\widehat\rho P_{\rm ready}),
 \\
 \mathcal I_{\rm fail}^t(\widehat\rho)
     =P_{\rm fail}\widehat\rho P_{\rm fail},
\end{gathered}
\label{eq:chvi-memory-complete-instrument}
\end{equation}
with multiple failure flags kept as separate blocks when supplied. Their
traces sum to one. Failure is a recorded outcome, not a missing attempt.

For the closed local protocol let the receiving region contain the memory
worldtube and read apparatus. Its access channel is the local channel just
specified, available after the finite write/copy completion. Retention is
tested for one read at any chosen time in its declared interval. Repeated
reads require their additional copy gates, blank qubits, disturbance and
work accounts; single-read performance does not automatically pay for them.

For a remote region, an additional causal delivery channel is an explicit
operand. Before its first allowed arrival $t_{\rm arr}$ the receiver output
has a common no-access flag, described by the trace-preserving replacement
map $\omega\mapsto\operatorname{Tr}(\omega)|\bot\rangle\langle\bot|$.
After arrival, a supplied delivery channel and its error can be composed
with the memory readout. The actual Chapter~IV hyperbolic propagation
constrains source-generated changes outside their causal future; neither
this memory coupling nor the word ``access'' supplies an instantaneous
remote control. The local construction is complete without selecting a
new remote carrier. Its remote extension must supply arrival and delivery
accuracy, rather than infer them from a finite register circuit.
Memory survival on $[0,T_{\rm store}]$ and remote readability after
$t_{\rm arr}$ are therefore distinct intervals.

\subsection{The resulting uniform error and resource ledger}

Let $\delta_{\rm acc}=\delta_{\rm read}+\epsilon_{\rm ap}$ include
the physical aperture term when that response is used. Let $\delta_{\rm init}$
bound any additional joint ready-state error for
the analog-pointer and blank-memory receiving preparation, relative to
the family for which the moment and momentum bounds above hold. Propagate
this error as a trace distance through the complete instrument. It does
not imply a second-moment bound for an arbitrary nearby actual state,
and it is not a second charge for marker preparation. First apply the
threshold and finite-write estimates to the finite-energy reference pipeline.
The actual marker preparation is within $\delta_{\rm prep}$ of that
reference at the original input; contraction through the entire actual
receiving channel adds $\delta_{\rm prep}$ once at its final output.
For compact notation write
$e_{\rm dec}^{\rm r}=e_{\rm dec,ref}^{\rm r}+\delta_{\rm prep}$,
as in the inherited decoder budget. This notation includes the final
whole-channel preparation payment; it does not transfer a moment bound
to an arbitrary nearby actual state.
Combining the preceding estimates gives the ready-conditional local bound
\begin{equation}
 \boxed{\begin{gathered}
 e_{\rm rec}^{\rm r}
 \leq\min\left\{1,
 e_{\rm dec}^{\rm r}
 +\frac{\mathcal B_{\rm ret}}{(a-w)^2}
 +r_w+\delta_{\rm init}+\Lambda_{T_{\rm store}}+\delta_{\rm acc}
 \right\}.
 \end{gathered}}
\label{eq:chvi-memory-ready-budget}
\end{equation}
All terms are per-label uniform bounds on the states where they are used.
The threshold argument handles bit/field correlations directly, and the
channel/hold arguments retain their environment. No independence between
decoding, analog noise, writing, storage and access is required.
For a complete attempt with failure probability at most $f$, counted as
unsuccessful, the bound becomes
\begin{equation}
 e_{\rm rec}^{\rm all}\leq f+(1-f)e_{\rm rec}^{\rm r}.
\label{eq:chvi-memory-all-attempt-budget}
\end{equation}
If failure was already included in an all-attempt channel error, it must
not be charged again by Eq.~\eqref{eq:chvi-memory-all-attempt-budget}.

A rectangular write has peak $\pi/\tau_w$, operator-norm action $\pi$,
and switch-work variation at most $2\pi/\tau_w$ in the common model chart.
Indeed its turn-on work lies in $[0,\pi/\tau_w]$ and its turn-off work in
$[-\pi/\tau_w,0]$. For finite-rest-energy inputs the signed total equals
the change of pointer rest
energy when the controller is zero at both endpoints; the spectator
source--field energy is conserved by this coupling. This signed work is
not identified with heat. The finite copy pulse has the same action and,
at duration $\tau_r$, peak $\pi/\tau_r$ and switch-work variation at most
$2\pi/\tau_r$. Its ideal diagonal-hold commuting limit has cancelling
signed on/off work; peak and temporary borrow remain charged.

The ledger additionally retains the upstream field-gate and bit-loading
translation controls, the pointer kinetic preparation, one fresh blank
memory bit and one read bit per supplied protocol, all failed preparation
attempts and their ready/failure
preparation accounts, any hold bath and coupling energy, physical delivery
when used, and later reset/disposal. A degenerate rest energy is not a
claim of free pure-state preparation or a Landauer heat formula. A cyclic
erasure bound would require its own reservoir, temperature and logical
ensemble. Model energies/work convert through the already declared
action/time units; no SI calibration is inferred.

For the proposed arithmetic recipient one may use
\begin{equation}
\begin{gathered}
 a=2^{10},\quad w=a/2=2^9,\quad k=2^{-40},
 \\ T_{\rm analog}=2^{20},\quad\tau_w=2^{-20},\\
 \sigma^2=C_{\rm prop}^2V_{\rm up}/2+kT_{\rm analog},\\
 p_2=\frac1{4\sigma^2},\\
 \mathcal B_{\rm ret}=\mathcal B_{\rm new}+2kT_{\rm analog}.
\end{gathered}
\label{eq:chvi-memory-evaluation-inputs}
\end{equation}
Here the propagation and detector ceilings retain their previously defined
source meanings; these equations specify the receiving pointer rather than
a measured calibration. The analog contribution is
$4\mathcal B_{\rm ret}/a^2$, and the finite-write allowance $r_w\leq10^{-4}$
is used only after checking Eq.~\eqref{eq:chvi-memory-write-discrepancy}
with the corresponding outward inputs.
A logical hold interval $T_{\rm store}=2^{20}$ with
$\Lambda_{T_{\rm store}}\leq10^{-4}$, together with
$\delta_{\rm init}\leq10^{-4}$ and $\delta_{\rm acc}\leq10^{-4}$,
is an explicit receiving specification, not an observed thermal lifetime.
The exact local copy/read reference can take $\delta_{\rm read}=0$;
a finite aperture still retains its separately bounded tail. For the example,
$\epsilon_{\rm ap}\leq2\cdot10^{-6}$ and
$\delta_{\rm read}\leq98\cdot10^{-6}$ fit the combined
$\delta_{\rm acc}\leq10^{-4}$ allowance.
The diagonal passive hold itself has $\Lambda_T=0$; a nonzero exposure
budget must be established from the actual admitted transverse interaction.
These distinctions keep a constructive finite record law, its mathematical
robustness, and achieved apparatus calibration as separately identified
parts of the same end-to-end account.

